% SPDX-License-Identifier: GPL-3.0-or-later
% Continuous layout exposes overlays for editing; Paged layout shows slides.
% Ctrl+R opens the compiled reference presentation.
\documentclass[aspectratio=169]{beamer}
\setbeamertemplate{navigation symbols}{}
\setbeamercolor{alerted text}{fg=red!75!black}
\title{Editing presentations with PasTeX}
\author{PasTeX contributors}
\date{}
\begin{document}

\begin{frame}[plain]
\titlepage
\end{frame}

\section{Editing slides}
\begin{frame}{Text and mathematics}{An editable widescreen frame}
Click a word or a formula to edit it through the native preview.
\[
  \int_0^1 x^2\,dx=\frac{1}{3},
  \qquad A=\begin{pmatrix}1&2\\3&4\end{pmatrix}.
\]
Use Fit Width or Fit Page to inspect the slide at different scales.
\end{frame}

\section{Overlays}
\begin{frame}{Revealing content}
This sentence appears on every overlay.
\begin{itemize}[<+->]
\item The first item appears immediately.
\item The second item appears on overlay two.
\item \alert{The last item is emphasized.}
\end{itemize}
\only<2>{This sentence occupies space only on overlay two.}
\uncover<3->{This sentence keeps its space while covered.}
\end{frame}

\begin{frame}{Pauses and alternatives}
Start here.\pause
\par Continue on the next overlay.\pause
\par \alt<3>{This is the third overlay.}{A different overlay is active.}
\par An alert can be selective: \alert<3>{watch this value}.
\[
  f(x)=\alert<3>{x^2}+\only<2>{1}\only<3->{2}.
\]
\end{frame}

\section{Layout}
\begin{frame}[t]{Columns and blocks}
\begin{columns}[T,onlytextwidth]
\begin{column}{.47\textwidth}
\begin{block}{A statement}
For every real number $x$, $x^2\geq0$.
\end{block}
\begin{exampleblock}{An example}
Taking $x=2$ gives $x^2=4$.
\end{exampleblock}
\end{column}
\begin{column}{.47\textwidth}
\begin{alertblock}{A restriction}
The ordering is an ordering of real numbers.
\end{alertblock}
Ordinary text and equations remain editable in each column.
\end{column}
\end{columns}
\end{frame}

\begin{frame}[t]{Replacing a statement}
The alternatives below occupy the same area in Paged layout.
Continuous layout exposes both alternatives for editing.
\begin{overprint}
\onslide<1|handout:1>
\begin{definition}[Nonnegative]
A real number $a$ is nonnegative if $a\geq0$.
\end{definition}
\onslide<2|handout:0>
\begin{theorem}[Squares]
For every real number $x$, the square $x^2$ is nonnegative.
\end{theorem}
\end{overprint}
\end{frame}

\begin{frame}[t]{Magnifying a region}
% The first overlay shows the whole frame; the second magnifies this region.
\framezoom<1><2>(0cm,0cm)(5cm,3cm)
This part is enlarged on the next overlay.
\[
  e^x=\sum_{n=0}^{\infty}\frac{x^n}{n!}.
\]
\end{frame}

\frame{\frametitle{Two frame syntaxes}
Frames can use either the frame environment or the frame command.
\par In the preamble, try
\texttt{\textbackslash usecolortheme\{dolphin\}} to change colors.
\par Native Beamer support is under development; use the compiled PDF to
check theme templates and advanced presentation features.
}
\end{document}
